% !PRINTED 510
\weberpage{510}
%<<CONT
be present. The totality $R$ of all common permutations of $Q$
and $Q'$ we call the \gesperrt{greatest common divisor},
or also the \gesperrt{intersection}, of $Q$ and $Q'$. This $R$ is always a
group, for if $\pi_1$ and $\pi_2$ both occur in $Q$ as well as in $Q'$,
then $\pi_1 \pi_2$ must also occur in $Q$ and in $Q'$, and hence also in $R$.
The concept carries over at once to several
groups $Q, Q', Q'' \ldots$

If $\psi$ is a function belonging to $Q$ and $\psi'$ one belonging to $Q'$,
we can determine the rational numbers $x, x'$ (by §. 143, 1.)
so that $\omega = x\psi + x'\psi'$ becomes a function belonging to $R$.
For $\omega$ certainly admits the permutations of $R$, since
$\psi$ and $\psi'$ admit them. If then $\pi$ is a permutation of $P$ not
contained in $R$, then certainly $\psi = \psi(\pi)$ and
$\psi' = \psi'(\pi)$ do not both hold; hence we can determine $x, x'$ so that
$\omega = \omega(\pi)$ likewise does not hold. If therefore we adjoin $\psi$ and $\psi'$,
and consequently $\omega$, simultaneously, the group of the field $N$
reduces by 3. to $R$. In the same way we may conclude for more than two groups $Q$,
and obtain the theorem:

4. \gesperrt{If $Q, Q', Q'' \ldots$ are divisors of $P$, $R$ their intersection,
$\psi, \psi', \psi'' \ldots$ functions belonging to $Q, Q', Q'' \ldots$,
then by the simultaneous adjunction of
$\psi, \psi', \psi'' \ldots$ the group of the
field $N$ reduces to $R$.}

If we adjoin not merely one root $\psi$ of the auxiliary equation
$\varphi(t) = 0$, but all the quantities conjugate to $\psi$,
namely $\psi, \psi_1, \psi_2 \ldots \psi_{j-1}$ — if, then, we derive from $\Omega$ the field

\[
\Omega'' = \Omega(\psi, \psi_1, \psi_2 \ldots \psi_{j-1})
\]

— then by this theorem the outcome is that the
group of $N$ in $\Omega''$ becomes the greatest common divisor $R$
of all the groups conjugate to $Q$, so that (by §. 154)
$R$ is the greatest common divisor of all the groups $\pi^{-1}Q\pi$,
as $\pi$ runs through the permutations of $P$.

\gesperrt{Of special interest now is the case in which all the conjugate
groups $\pi^{-1}Q\pi$ are identical with one another.} In this
case, by Theorem 2., each of the conjugate quantities
$\psi, \psi_1, \psi_2 \ldots \psi_{j-1}$ is contained in $\Omega(\psi)$, the fields
$\Omega(\psi), \Omega(\psi_1) \ldots \Omega(\psi_{j-1})$ and $\Omega(\psi, \psi_1 \ldots \psi_{j-1})$ are identical,
and $\Omega(\psi)$ is a \gesperrt{normal field} over $\Omega$.
